% !TEX program = lualatex
% Compile from 2026/week4; see README-beamer.md for the merged PDF steps (main output: .beamer-build/main).
% Source: ../../2025/week4/week4.pdf (19 slides, 4:3) and week4.ipynb.
\documentclass[12pt,aspectratio=43]{beamer}
% Shared Beamer design. Load after \documentclass using \input.
% Engine: LuaLaTeX. Fonts: UD Digi Kyokasho N (regular/bold), Arial.
\usepackage{luatexja-fontspec}
\setmainjfont[BoldFont={UD Digi Kyokasho N Bold}]{UD Digi Kyokasho N}
\setsansjfont[BoldFont={UD Digi Kyokasho N Bold}]{UD Digi Kyokasho N}
\setmonojfont[BoldFont={UD Digi Kyokasho N Bold}]{UD Digi Kyokasho N}
% Liberation Sans keeps the layout usable on systems without Arial.
\IfFontExistsTF{Arial}{
  \setmainfont{Arial}
  \setsansfont{Arial}
  \setmonofont{Arial}
}{
  \setmainfont{Liberation Sans}
  \setsansfont{Liberation Sans}
  \setmonofont{Liberation Sans}
}
\renewcommand{\kanjifamilydefault}{\gtdefault}

% Typeset mathematics with a dedicated TeX math font, independent of Arial.
\usefonttheme{professionalfonts}
\usepackage[mathrm=sym]{unicode-math}
\setmathfont{Latin Modern Math}
\setoperatorfont{\symup}

\definecolor{courseblue}{HTML}{4472C4}
\definecolor{ruleblue}{HTML}{8FAADC}
\setbeamercolor{normal text}{fg=black,bg=white}
\setbeamercolor{structure}{fg=courseblue}
\setbeamercolor{frametitle}{fg=black,bg=white}
\setbeamercolor{title banner}{fg=white,bg=courseblue}
\setbeamerfont{frametitle}{size=\fontsize{18}{22}\selectfont,series=\mdseries}
\setbeamerfont{quote}{shape=\upshape}
\setbeamersize{text margin left=9mm,text margin right=9mm}
\setbeamertemplate{navigation symbols}{}
\setbeamercolor{page number}{fg=gray,bg=white}
\setbeamertemplate{footline}{%
  \begin{beamercolorbox}[wd=\paperwidth,ht=2.5mm,dp=1.5mm]{page number}
    \hfill{\fontsize{8}{9}\selectfont
      \ifnum\value{framenumber}>0
        \insertframenumber\,/\,\inserttotalframenumber
      \fi}\hspace*{3mm}
  \end{beamercolorbox}}
\setbeamertemplate{headline}{%
  \hfill\textcolor{gray}{\fontsize{6}{7}\selectfont
    \hypersetup{linkcolor=gray}% Keep Beamer's title links in the header color.
    \ifnum\value{framenumber}=0 \insertshortdate\enspace\fi
    \insertshorttitle\space\insertshortsubtitle}\hspace{2mm}\vspace{1.5mm}}
\setbeamertemplate{frametitle}{%
  \vspace{2mm}\centering\insertframetitle\par
  \vspace{-1mm}{\color{ruleblue}\rule{.44\paperwidth}{.5pt}}\par
  \vspace{2mm}}
\setbeamercolor{itemize item}{fg=courseblue}
\setbeamercolor{itemize subitem}{fg=black}
\setbeamercolor{itemize subsubitem}{fg=black}
\setbeamertemplate{itemize item}{\textbullet}
\setbeamertemplate{itemize subitem}{\textnormal{-}}
\setbeamertemplate{itemize subsubitem}{\textnormal{-}}
\setlength{\parskip}{.45em}
\hypersetup{colorlinks=true,urlcolor=courseblue}

% Use \titlepage inside a [t,noframenumbering] frame. The background extends
% 1 mm past both page edges to avoid hairline gaps in PDF viewers.
\setbeamertemplate{title page}{%
  \vspace{2mm}
  \begin{beamercolorbox}[wd=\paperwidth,ht=43mm,dp=13mm,colsep*=1mm,center]{title banner}
    {\fontsize{32}{40}\selectfont\inserttitle\par}
    \vspace{2mm}
    {\fontsize{20}{24}\selectfont\insertsubtitle\par}
  \end{beamercolorbox}
  \vspace{10mm}
  \hfill\insertauthor
}

\usepackage{listings}
\usepackage{upquote}
\IfFontExistsTF{Consolas}{\newfontfamily\codefont{Consolas}}{\newfontfamily\codefont{Liberation Mono}}
\definecolor{codevalue}{HTML}{7A3E9D}
\definecolor{codebuiltin}{HTML}{006DB0}
\definecolor{codestring}{HTML}{267348}
\newcommand{\code}[1]{{\codefont #1}}

\lstdefinestyle{python}{
  language=Python,
  basicstyle=\codefont\fontsize{14}{16}\selectfont,
  keywordstyle=\color{codebuiltin},
  keywordstyle=[2]\color{codebuiltin},
  morekeywords=[2]{print,range},
  stringstyle=\color{codestring},
  showstringspaces=false,
  columns=fullflexible,
  keepspaces=true,
  upquote=true,
  tabsize=4,
  breaklines=false,
  backgroundcolor=\color{black!4},
  frame=single,framerule=0pt,framesep=1.7mm,
  xleftmargin=1.7mm,xrightmargin=1.7mm,
  aboveskip=1mm,belowskip=0.5mm,
  literate=*{0}{{{\color{codevalue}0}}}1
    {1}{{{\color{codevalue}1}}}1 {2}{{{\color{codevalue}2}}}1
    {3}{{{\color{codevalue}3}}}1 {4}{{{\color{codevalue}4}}}1
    {5}{{{\color{codevalue}5}}}1 {6}{{{\color{codevalue}6}}}1
    {7}{{{\color{codevalue}7}}}1 {8}{{{\color{codevalue}8}}}1
    {9}{{{\color{codevalue}9}}}1
}
\lstset{style=python}
\newcommand{\result}[1]{%
  \par\vspace{1mm}%
  {\small\textcolor{black!60}{出力}}\quad
  {\codefont\fontsize{16}{18}\selectfont #1}\par}
\setbeamercolor{lesson note}{fg=black,bg=courseblue!9}
\newenvironment{lessonbox}
  {\begin{beamercolorbox}[wd=\linewidth,sep=1.5mm]{lesson note}}
  {\end{beamercolorbox}\par}

\title{計算物理学B}
\subtitle{第四回：配列（NumPy）}
\author{野垣 康介、藤本 悠輝}
\date{2026/11/10}
\hypersetup{pdftitle={計算物理学B 第四回：配列（NumPy）},pdfauthor={野垣 康介、藤本 悠輝}}

\begin{document}

% 0. 表紙
\begin{frame}[t,noframenumbering]
  \titlepage
\end{frame}

% 1. 2026年度の講義予定
\begin{frame}{講義予定}
  \setlength{\parskip}{0pt}
  2026年度 第2学期（火曜3限）
  \par\vspace{2mm}
  \begingroup\centering
    \renewcommand{\arraystretch}{1.05}
    \begin{tabular}{@{}l@{\hspace{2mm}}l@{\hspace{4mm}}l@{\hspace{2mm}}l@{}}
      10/6 & 四則演算 & 12/8 & モンテカルロ1 \\
      10/27 & 制御文（for, if） & 12/15 & モンテカルロ2 \\
      11/6（金） & 関数 & 12/22 & 微分方程式1 \\
      \textcolor{courseblue}{\textbf{11/10}} &
      \textcolor{courseblue}{\textbf{配列（NumPy）}} & 1/12 & 微分方程式2 \\
      11/17 & 可視化（matplotlib） & 1/19 & 偏微分方程式1 \\
      11/24 & 数値微分 & 1/26 & 偏微分方程式2 \\
      12/1 & 数値積分 & 2/2 & 根の探索 \\
      \multicolumn{2}{c}{\small 〜中間レポート〜} &
      \multicolumn{2}{c}{\small 〜期末レポート〜}
    \end{tabular}\par
  \endgroup
  \vspace{2mm}
  {\small
    \textcolor{courseblue}{10/20は出張で休講。}\\
    11/6（金）は火曜授業の振替。\par}
  \vspace{1mm}
  {\fontsize{7}{8}\selectfont
    日付：\href{https://www.niigata-u.ac.jp/wp-content/uploads/2025/08/schedule_2026.pdf}{令和8年度（2026）新潟大学授業暦}\par}
\end{frame}

% 2. 授業で使うリンク
\begin{frame}{授業で使うリンク}
  \begin{itemize}
    \setlength{\itemsep}{4mm}
    \item Google Colab\\
      {\small\url{https://colab.research.google.com/?hl=ja}}
    \item 講義資料（GitHub）\\
      {\small\url{https://github.com/nogaki/Computational_Physics_B}}
    \item 今週の教材：\code{2026/week4/week4.ipynb}
  \end{itemize}
\end{frame}

% 3. GitHubからColabへ取り込む
\begin{frame}[t]{GitHubからColabへ取り込む}
  \setlength{\parskip}{0pt}
  \begin{enumerate}
    \setlength{\itemsep}{4mm}
    \item GitHubで\code{week4.ipynb}を開き、\\
      ブラウザーのアドレス欄のURLをコピーします。
    \item Colabで「ファイル」→「ノートブックを開く」\\
      →「GitHub」を選びます。
    \item コピーしたURLを貼り付けて検索し、表示されたノートブックを開きます。
  \end{enumerate}
  \vspace{3mm}
  \begin{lessonbox}
    \textbf{開いたら「ドライブにコピーを保存」}\\[1mm]
    \small 「ファイル」メニューから、自分用のコピーを保存します。\\
    開いただけでは、自分のドライブに保存されません。
  \end{lessonbox}
\end{frame}

% 4. 物理量を時刻ごとに並べる
\begin{frame}[t]{物理量を時刻ごとに並べる}
  \setlength{\parskip}{0pt}
  \small
  空気抵抗を無視し、初速0、落下開始位置\(x=0\)、下向きを正として自由落下を考えます。
  \par\vspace{1mm}
  \(g=10.0\ \mathrm{m/s^2}\)と近似すると、落下開始からの経過時間\(t\geq0\)で
  \(x=\frac{1}{2}gt^2\)、\(v=gt\)です。
  \par\vspace{2mm}
  \begingroup\centering
    \renewcommand{\arraystretch}{1.2}
    \begin{tabular}{@{}l@{\hspace{8mm}}rrrr@{}}
      時刻 \(t\,[\mathrm{s}]\) & 0 & 1 & 2 & 3 \\
      位置 \(x\,[\mathrm{m}]\) & 0 & 5 & 20 & 45 \\
      速度 \(v\,[\mathrm{m/s}]\) & 0 & 10 & 20 & 30
    \end{tabular}\par
  \endgroup
  \vspace{2mm}
  \begin{lessonbox}
    \small 各列は同じ時刻の物理量です。時刻の点数を増やしても、この対応を保って計算します。
  \end{lessonbox}
\end{frame}

% 5. 時刻と位置をリストにする
\begin{frame}[t,fragile]{時刻と位置をリストにする}
  \setlength{\parskip}{0pt}
  時刻と位置を同じ順番のリストに並べます。\code{times}の単位は秒、\code{positions}はメートルです。
\begin{lstlisting}
times = [0.0, 1.0, 2.0]
positions = [0.0, 5.0, 20.0]
print(times[1], positions[1])
print(len(times))
\end{lstlisting}
  \result{1.0\quad 5.0}
  \result{3}
  \begin{lessonbox}
    \small 同じ添字が同じ時刻を表します。添字は0から始まり、ここでは\code{times[1]}と\code{positions[1]}が1秒と5メートルです。
  \end{lessonbox}
\end{frame}

% 6. 計算結果を追加して取り出す
\begin{frame}[t,fragile]{計算結果を追加して取り出す}
  \setlength{\parskip}{0pt}
  3秒後の位置\(45.0\ \mathrm{m}\)も、時刻と位置の両方に追加します。
\begin{lstlisting}
times.append(3.0)
positions.append(45.0)
for i in range(len(times)):
    print(times[i], positions[i])
\end{lstlisting}
  \begin{columns}[T,onlytextwidth]
    \begin{column}{.28\textwidth}
      {\small\textcolor{black!60}{出力（各行）}}\\[-1mm]
      {\codefont\fontsize{14}{16}\selectfont
        0.0 0.0\\
        1.0 5.0\\
        2.0 20.0\\
        3.0 45.0\par}
    \end{column}
    \begin{column}{.69\textwidth}
      \small
      \textbf{リスト名の後に付ける操作}\\[1mm]
      \code{.append(x)}：末尾に\code{x}を追加\\[1mm]
      \code{.pop()}：末尾を削除して値を返す\\[1mm]
      \code{.pop(i)}：添字\code{i}の要素を削除して返す\\[1mm]
      \code{.remove(x)}：最初に一致する\code{x}を削除
    \end{column}
  \end{columns}
  \vspace{1mm}
  {\small 時刻と位置の両リストから、同じ添字の要素を削除します。}
\end{frame}

% 7. 前回の関数を使う
\begin{frame}[t,fragile]{前回の関数を使う}
  \setlength{\parskip}{0pt}
  前回の\code{free\_fall}で、3秒後の位置と速度を求めます。
\begin{lstlisting}
def free_fall(t):
    g = 10.0
    x = 0.5 * g * t ** 2
    v = g * t
    return x, v

x_one, v_one = free_fall(3.0)
print(x_one, v_one)
\end{lstlisting}
  \result{45.0\quad 30.0}
  \small 出力は位置（m）、速度（m/s）の順です。
\end{frame}

% 8. 配列を作る関数と配列の型
\begin{frame}[t,fragile]{arrayとndarrayの違い}
  \setlength{\parskip}{0pt}
  \small \code{np.array}は配列を作る関数、\code{ndarray}はその型です。
\begin{lstlisting}
import numpy as np
t = np.array([0.0, 1.0, 2.0, 3.0])
print(type(t))
\end{lstlisting}
  \result{<class 'numpy.ndarray'>}
  \vspace{2mm}
  \begingroup
    \small
    \renewcommand{\arraystretch}{1.1}
    \begin{tabular}{@{}l@{\hspace{3mm}}p{7.0cm}@{}}
      \code{list} & \code{[]}で作る組込みのリスト。型を混在でき、\mbox{\code{* 2}}は繰り返し。\\
      \code{array.array} & Python標準ライブラリの1次元配列。型を揃えて保存し、\mbox{\code{* 2}}は繰り返し。\\
      \code{numpy.ndarray} & 1次元も含む多次元配列。要素型（\code{dtype}）を揃え、\mbox{\code{* 2}}は各要素を2倍。
    \end{tabular}\par
  \endgroup
\end{frame}

% 9. 配列を使った一括計算と時刻の点数
\begin{frame}[t,fragile]{配列で全時刻を計算する}
  \setlength{\parskip}{0pt}
  \small この関数の\code{*}と\code{**}は、配列の各要素に作用します。
\begin{lstlisting}
x, v = free_fall(t)
print(x)
print(v)
\end{lstlisting}
  {\small\textcolor{black!60}{出力}}\quad
  {\codefont\fontsize{12}{14}\selectfont
    \begin{tabular}{@{}ll@{}}
      x & [0. 5. 20. 45.]\\
      v & [0. 10. 20. 30.]
    \end{tabular}\par}
  \vspace{.5mm}
  \small 時刻の点数を増やすときも、同じ関数を使えます。
\begin{lstlisting}
t = np.linspace(0.0, 3.0, 7)
x, v = free_fall(t)
\end{lstlisting}
  \small 0秒と3秒を含む7点、間隔0.5秒です。\\
  \code{x[i]}と\code{v[i]}は、同じ時刻\code{t[i]}に対応します。
\end{frame}

% 10. 同じ時刻のデータを取り出す
\begin{frame}[t,fragile]{同じ時刻のデータを取り出す}
  \setlength{\parskip}{0pt}
\begin{lstlisting}
print(t[2], x[2], v[2])
print(t[:3])
print(x[:3])
\end{lstlisting}
  \result{1.0\quad 5.0\quad 10.0}
  \result{[0. 0.5 1.]}
  \result{[0. 1.25 5.]}
  \small 添字2は1秒で、位置5メートル・速度10メートル毎秒です。\code{[:3]}は添字0,1,2を取り出し、終点3は含みません。
\end{frame}

% 11. 各時刻の運動エネルギー
\begin{frame}[t,fragile]{各時刻の運動エネルギー}
  \setlength{\parskip}{0pt}
  質量\(m=2.0\ \mathrm{kg}\)の物体の運動エネルギーは\(E=\frac{1}{2}mv^2\)です。
\begin{lstlisting}
mass = 2.0
energy = 0.5 * mass * v ** 2
print(energy)
\end{lstlisting}
  {\small\textcolor{black!60}{出力}}\quad
  {\codefont\fontsize{12}{14}\selectfont [0. 25. 100. 225. 400. 625. 900.]\par}
  \small \code{energy[i]}は\code{t[i]}での運動エネルギー（単位\(\mathrm{J}\)）です。静止からの自由落下では、失った位置エネルギー\(mgx\)に等しくなります。
\end{frame}

% 12. 多次元配列の行・列・添字
\begin{frame}[t,fragile]{多次元配列：行と列}
  \setlength{\parskip}{0pt}
  \small 1次元は添字1つ、2次元は行・列の2つで要素を指定します。
  \begin{columns}[T,onlytextwidth]
    \begin{column}{.46\textwidth}
\begin{lstlisting}
r = np.array([
    [1, 0],
    [0, 2],
    [-1, 0]
])
\end{lstlisting}
    \end{column}
    \begin{column}{.51\textwidth}
      \small 行は粒子、列は座標（単位m）。\\[1mm]
      \begingroup
        \setlength{\tabcolsep}{3pt}
        \renewcommand{\arraystretch}{1.05}
        \begin{tabular}{@{}crr@{}}
          & \textcolor{courseblue}{列0（\(x\)）} & \textcolor{courseblue}{列1（\(y\)）}\\
          \hline
          行0 & 1 & 0\\
          行1 & 0 & 2\\
          行2 & -1 & 0
        \end{tabular}\par
      \endgroup
      \code{r.shape} → \code{(3, 2)}\\
      3行・2列の2次元配列です。
    \end{column}
  \end{columns}
  \vspace{1mm}
  \begingroup
    \small
    \setlength{\tabcolsep}{3pt}
    \renewcommand{\arraystretch}{1.1}
    \begin{tabular}{@{}lll@{}}
      \code{r[1, 0]} & \code{0} & 行1・列0の要素\\
      \code{r[1]} & \code{[0 2]} & 行1の全体（2番目の粒子）\\
      \code{r[:, 0]} & \code{[1 0 -1]} & 全行の列0（\code{:}は全て）
    \end{tabular}\par
  \endgroup
  \vspace{1mm}
  {\footnotesize 時刻ごとにこの表を重ねると、（時刻数, 3, 2）の3次元配列になります。}
\end{frame}

% 13. 行列で全粒子を回転する
\begin{frame}[t,fragile]{行列で全粒子を回転する}
  \setlength{\parskip}{0pt}
  粒子の位置を原点のまわりに反時計回り90度回転します。
  \[\begin{pmatrix}x'\\y'\end{pmatrix}=
    R\begin{pmatrix}x\\y\end{pmatrix},\qquad
    R=\begin{pmatrix}0&-1\\1&0\end{pmatrix}.\]
  \small 行データには右から\(R^T\)を掛けます。\code{@}は行列積、\code{.T}は転置、\code{*}は要素積です。
\begin{lstlisting}
R = np.array([[0, -1], [1, 0]])
rotated = r @ R.T
print(rotated)
\end{lstlisting}
  {\small\textcolor{black!60}{出力}}\quad
  {\codefont\fontsize{12}{14}\selectfont \(\begin{bmatrix}0&1\\-2&0\\0&-1\end{bmatrix}\)\par}
\end{frame}

% 14. 繰り返し測った周期のばらつき
\begin{frame}[t,fragile]{繰り返し測った周期のばらつき}
  \setlength{\parskip}{0pt}
  \small 仮想の振り子周期4件です（実測値ではなく、単位\(\mathrm{s}\)）。
\begin{lstlisting}
periods = np.array([1.75, 2.0,
                    2.25, 2.0])
print(np.mean(periods))
print(np.var(periods))
\end{lstlisting}
  {\small 出力：平均\(2.0\ \mathrm{s}\)、分散\(0.03125\ \mathrm{s^2}\)。}
  {\small 測定値\(T_i\)（\(i=1,\ldots,n\)）、測定回数\(n\)：}
  {\footnotesize
    \(\bar{T}=\frac{1}{n}\sum_{i=1}^{n}T_i\)\qquad
    \(\sigma_T^2=\frac{1}{n}\sum_{i=1}^{n}(T_i-\bar{T})^2\)\par}
  \par\vspace{1mm}
  {\small \code{np.var}の既定分母は\(n\)。分散は測定値の散らばりで、平均値の誤差とは別です。}
\end{frame}

% 15. 仮想的な周期測定データを作る
\begin{frame}[t,fragile]{仮想的な周期測定データを作る}
  \setlength{\parskip}{0pt}
  真の周期は\(T_0=2.0\ \mathrm{s}\)。誤差は独立な正規分布（平均0、標準偏差\(0.1\ \mathrm{s}\)）とする仮想モデルです。
\begin{lstlisting}[aboveskip=1mm,belowskip=0mm]
rng = np.random.default_rng(0)
samples = rng.standard_normal(1000)
samples = 2.0 + 0.1 * samples
values = samples.tolist()
\end{lstlisting}
  {\small 1000回分の仮想測定です。理論平均\(2.0\ \mathrm{s}\)、分散\(0.01\ \mathrm{s^2}\)ですが有限標本ではずれます。\\
    \code{values}は同じデータのリスト。seedを固定すると同じ環境で再現しやすくなります。}
\end{frame}

% 16. 物理の例で配列を使う
\begin{frame}[t]{物理の例で配列を使う}
  \setlength{\parskip}{0pt}
  \small
  \begin{enumerate}
    \setlength{\itemsep}{3mm}
    \item \(0\)〜\(3\)秒の31点で、自由落下の位置・速度と、質量\(2.0\ \mathrm{kg}\)の運動エネルギーを配列で求めましょう。
    \item 3粒子を90度回転した前後で、原点からの距離の2乗が変わらないか確かめましょう。ヒント：\code{r[:, 0]**2 + r[:, 1]**2}。
    \item \code{samples}の平均と分散を求め、理論値の\(2.0\ \mathrm{s}\)、\(0.01\ \mathrm{s^2}\)と比べましょう。
  \end{enumerate}
  \vfill
  {\small \code{for}文による統計計算と実行時間の比較は、ノートブックの補足・発展課題にあります。}
\end{frame}

\end{document}
